



\section{Mã hóa bằng tay thuật toán AES}

PLAINTEXT = \textbf{DLK24DHGTVTTPHCM} \\
CIPHER KEY =\textbf{LEHUYNHANHKHOI24}

\subsection{Quy trình mở rộng khóa phụ}

Thực hiện thuật toán mở rộng khóa nhằm sinh khóa phụ cho các vòng lặp mã hóa (\textbf{KeyExpension}) như sau:

\begin{enumerate}[label= \textbf{Bước \arabic*:}]
	\item Điền $N_k = 4$ từ đầu tiên bằng khóa mã hóa. Mỗi từ $W[i]$ tiếp theo sẽ được điền bằng kết quả phép \textbf{XOR} giữa từ $W[i-1]$ với $W[i-N_k]$ ($i \geq 4$). Đối với những trường hợp từ $W[i]$, với $i$ là bội số của $N_k$, thì ta cần thực hiện $transW(W[i-1])$, bao gồm các phép chuyển đổi $RWORD$, $SubBytes$ và $RCON$ với $W[i-1]$ rồi sau đó thực hiện phép $XOR$ với $W[i-N_k]$
	
	\textbf{Ví dụ:}  Ta thực hiện mở rộng khóa, bắt đầu bằng $W[4]$:
	\begin{itemize}
		\item Thực hiện phép $RWORD$, tức là biến từ đầu vào $(a,b,c,d) \text{ thành đầu ra }(b,c,d,a).$
		\begin{center}
			$
			W[3] = 	\begin{bmatrix} 4F \\49 \\32 \\34 \end{bmatrix}
			\xrightarrow{ROTWORD}
			\begin{bmatrix}	49 \\32 \\34 \\4F \end{bmatrix}
			$
		\end{center}
		
		\item Thực hiện phép $SubBytes$ với bảng tham chiếu \textbf{S-box}:
		\begin{center}
			$
			W[3] = \begin{bmatrix}	49 \\32 \\34 \\4F \end{bmatrix}
			\xrightarrow{SUBBYTES}
			\begin{bmatrix}	3B \\ 23 \\18\\84 \end{bmatrix}
			$
		\end{center}
		
		\item Thực hiện phép $RCON$, tức là ta sẽ XOR kết quả của phép \textbf{SubBytes} với hằng số vòng $RCON[i]$:
		\begin{center}
			$transW(W[3]) = W[4] \oplus RCON[1] = \begin{bmatrix}	3B \\ 23 \\18\\84 \end{bmatrix} \oplus \begin{bmatrix}	01 \\ 00 \\00\\00 \end{bmatrix}
			= \begin{bmatrix} 3A\\23\\18\\84\end{bmatrix}$ 
		\end{center}
		\item Kết quả của phép $RCON$ cũng chính là đầu ra của phép biến đổi từ $W[i-1]$ với $i$ là bội số của $N_k$. Sau đó, ta thực hiện phép $XOR$ giữa kết quả phép biến đổi $W[i-1]$ với $W[i-N_k]$. Đầu ra của phép XOR này chính là từ $W[i]$ tiếp theo. 
		\begin{center}
			$W[4] = transW(W[3]) \oplus W[0] = \begin{bmatrix} 3A\\23\\18\\84\end{bmatrix} \oplus \begin{bmatrix} 4C\\45\\48\\55\end{bmatrix} = \begin{bmatrix} 76\\66\\50\\D1\end{bmatrix}$ 
		\end{center}
	\end{itemize}
	
	\item Tương tự với $W[5], W[6], W[7]$ ta thu được:
	
	\begin{center}
		$W[5] = W[4] \oplus W[1] = \begin{bmatrix} 76\\66\\50\\D1\end{bmatrix} \oplus \begin{bmatrix} 4C\\45\\48\\55\end{bmatrix} = \begin{bmatrix} 2F\\28\\18\\90\end{bmatrix}$\\
		$W[6] = W[5] \oplus W[2] = \begin{bmatrix} 2F\\28\\18\\90\end{bmatrix} \oplus \begin{bmatrix} 59\\4E\\48\\41\end{bmatrix} =  \begin{bmatrix} 61\\60\\53\\D8\end{bmatrix}$\\
		$W[7] = W[6] \oplus W[3] = \begin{bmatrix} 61\\60\\53\\D8\end{bmatrix} \oplus \begin{bmatrix} 4F\\49\\32\\34\end{bmatrix} = \begin{bmatrix} 2E\\29\\61\\EC\end{bmatrix}$\\
	\end{center}
	
	\item Sau khi điền $W[7]$, ta thu được khóa phụ thứ nhất ($RoundKey1= [W[4],W[5],W[6],W[7]]$):
	\begin{center}
		$
		RoundKey1 = 
		\begin{bmatrix}
			76 & 2F & 61 & 2E\\
			66 & 28 & 60 & 29 \\
			50 & 18 & 53 & 61 \\
			D1 & 90 & D8 & EC
		\end{bmatrix}
		$
	\end{center}
	
	\item Thực hiện sinh khóa các vòng sau tương tự các bước trên, cho đến khi ta thu được khóa phụ cho vòng lặp mã hóa thứ 10 (\textbf{RoundKey10}).
	
\end{enumerate}

\subsection{Quy trình mã hóa}
Vì cả chiều dài của bản khối và khóa đều tương ứng với 128bit $\implies N_b = 4$ \& $N_k = 4$
$\implies$ AES sẽ được thực hiện với số vòng lặp là 10.


Quy trình thực hiện thuật toán mã hóa đối xứng AES như sau: 

\begin{enumerate}[label= \textbf{Bước \arabic*:}]
	\item Lần lượt chuyển đổi khối dữ liệu bản rõ $128 bit$ được chia thành $16 bytes$, tạo thành Ma trận Trạng thái (State) với kích cỡ $4x4$, tương tự với khóa mã hóa $128 bit$.
	
	% 44 4C 4B 32 34 44 48 47 54 56 54 54 50 48 43 4D (DLK24DHGTVTTPHCM)
	% 4C 45 48 55 59 4E 48 41 4E 48 4B 48 4F 49 32 34 (LEHUYNHANHKHOI24)
	
	
	\begin{center}
		\begin{tabular}{cc}
			$\begin{bmatrix}
				44 & 34 & 54 & 50 \\
				4C & 44 & 56 & 48 \\
				4B & 48 & 54 & 43 \\
				32 & 47 & 54 & 4D
			\end{bmatrix}$ 
			& 
			$\begin{bmatrix}
				4C & 59 & 4E & 4F \\
				45 & 4E & 48 & 49 \\
				48 & 48 & 4B & 32 \\
				55 & 41 & 48 & 34
			\end{bmatrix}$ \\
			(a) Ma trận State & (b) Ma trận Cypher Key
		\end{tabular}
	\end{center}
	
	\item Thực hiện hàm \textbf{AddRoundKey}, nghĩa là thực hiện $State \oplus Round Key$ với mỗi bytes của States tại vị trí $[i][j]$ và bytes của RoundKey tại vị trí $[i][j]$ tương ứng. 
	
	Tại bước khởi tạo \textbf{AddRoundKey} đầu tiên, ta sẽ khởi tạo bằng cách thực hiện phép \textbf{XOR} giữa \textbf{Plaintext} với \textbf{CypherKey} ban đầu. 
	
	Hàm \textbf{AddRoundKey} ở các vòng tiếp theo, ta sẽ thực hiện phép \textbf{XOR} giữa output các phép biến đổi ở vòng trước kết hợp với khóa phụ \textbf{RoundKey} tương ứng.
	
	\textbf{Ví dụ:} Tại vị trí $[1][0]$,  $44 \oplus 4C = 0100 0100 \oplus 0100 1100 =0000 1000 = 8$
	
	\begin{align*}
		\text{AddRoundKey} &= \text{StateMatrix} \oplus \text{CypherKey} \\
		\text{AddRoundKey} &= 
		\begin{bmatrix}
			44 & 48 & 54 & 50 \\
			4C & 47 & 56 & 48 \\
			32 & 54 & 54 & 43 \\
			44 & 56 & 54 & 4D
		\end{bmatrix}
		\oplus
		\begin{bmatrix}
			4C & 59 & 4E & 4F \\
			45 & 4E & 48 & 49 \\
			48 & 48 & 4B & 32 \\
			55 & 41 & 48 & 34
		\end{bmatrix} \\
		\text{AddRoundKey} &= 
		\begin{bmatrix}
			08 & 6d & 1A & 1F \\
			09 & 0A & 1E & 01 \\
			03 & 00 & 1F & 71 \\
			67 & 06 & 1C & 79
		\end{bmatrix}
	\end{align*}
	
	
	%SubBytes%
	\item Thực hiện phép thay phi tuyến tính thông qua hàm \textbf{SubBytes}, tức là ta sẽ thay thế từng byte trong output của hàm \textbf{AddRoundKey} bằng một bảng quy định giá trị thay thế, S-box \ref{tab:sbox}:
	
	\textbf{Ví dụ}: Tại vị trí $[1][0]$, byte tại vị trí này là ${08}$, ứng với bảng \textbf{S-box} tại hàng số 0, cột số 8 sẽ là byte $30$
	\begin{center}
		$\begin{bmatrix}
			08 & 6d & 1A & 1F \\
			09 & 0A & 1E & 01 \\
			03 & 00 & 1F & 71 \\
			67 & 06 & 1C & 79
		\end{bmatrix}
		\xrightarrow{SubBytes}
		\begin{bmatrix}
			30 & 3C & A2 & C0\\
			01 & 67 & 72 & 7C\\
			7B & 63 & C0 & A3\\
			85 & 6F & 9C & B6 \\
		\end{bmatrix}
		$
	\end{center}
	
	%ShiftRows%
	\item \textbf{ShiftRows} có chức năng quay trái từng hàng trạng thái với hệ số quay $\{C_1, C_2, C_3\}$. Hàng đầu tiên (hàng số 0) sẽ không dịch chuyển. Tùy thuộc vào giá trị của $N_b$ mà hệ số quay sẽ tương ứng với chúng. Với $N_b = 4$: 
	\begin{itemize}
		\item Hàng 1 được dịch chuyển bởi $C_1 = 1$ byte
		\item Hàng 2 được dịch chuyển bởi $C_2 = 2$ byte
		\item Hàng 3 được dịch chuyển bởi $C_3 = 3$ byte
	\end{itemize}
	\begin{center}
		$
		\begin{bmatrix}
			30 & 3C & A2 & C0\\
			01 & 67 & 72 & 7C\\
			7B & 63 & C0 & A3\\
			85 & 6F & 9C & B6 \\
		\end{bmatrix}
		\xrightarrow{ShiftRows}
		\begin{bmatrix}
			30 & 3C & A2 & C0\\
			67 & 72 & 7C & 01\\
			C0 & A3 & 7B & 63\\
			B6 & 85 & 6F & 9C\\
		\end{bmatrix}
		$
		
	\end{center}
	
	\item \textbf Phép biến đổi {MixColumns}, thực hiện tính toán từng cột của state. Các cột trong state được coi như một đa thức trong trường $GF(2^8)$. Việc biến đổi cột trạng thái được thực hiện vởi phép nhân ($\times$) và phép $\text{XOR }(\oplus)$.
	
	\textbf{Ví dụ:} Ta sẽ biến đổi cột đầu tiên $c_0$ của output \textbf{ShiftRows} bằng cách thực hiện phép biến đổi \textbf{MixColumns}:
	\begin{center}
		$c_0
		=
		\begin{bmatrix}
			02 & 03 & 01 & 01 \\
			01 & 02 & 03 & 01 \\
			01 & 01 & 02 & 03\\
			03 & 01 & 01 & 02 \\
		\end{bmatrix} \begin{bmatrix} 30 \\ 67 \\ C0 \\B6 \end{bmatrix}
		=  \begin{bmatrix} BF \\ 13 \\ 0D \\ 80 \end{bmatrix}$
	\end{center}
	
	Xét phần tử đầu tiên của $c_0$: 
	\begin{align*}
		c_0[0] &= (02 \cdot 30) \oplus (03 \cdot 67) \oplus (01 \cdot C0) \oplus (01 \cdot B6)\\
		c_0[0] &= 60 \oplus A9 \oplus C0 \oplus B6\\
		c_0[0] &= BF
	\end{align*}
	
	Thực hiện tương tự với các cột còn lại, ta thu được ouput của bước \textbf{MixColumns} như sau:
	\begin{center}
		$
		\begin{bmatrix}
			30 & 3C & A2 & C0\\
			67 & 72 & 7C & 01\\
			C0 & A3 & 7B & 63\\
			B6 & 85 & 6F & 9C\\
		\end{bmatrix}
		\xrightarrow{MixColumns}
		\begin{bmatrix}
			BF & C8 & CF & 67\\
			13 & A3 & B8 & FB\\
			0D & 87 & 99 & B8\\
			80 & 84 & 24 & 1A\\
		\end{bmatrix}
		$ 
		
		
	\end{center}
	
	\item Thực hiện hàm \textbf{AddRoundKey} nhằm thêm khóa phụ cho vòng lặp thứ hai, ta thực hiện XOR giữa State (đầu ra của bước \textbf{MixColumns}) với \textbf{RoundKey1}
	\begin{center}
		$
		\begin{bmatrix}
			BF & C8 & CF & 67\\
			13 & A3 & B8 & FB\\
			0D & 87 & 99 & B8\\
			80 & 84 & 24 & 1A\\
		\end{bmatrix}
		\oplus
		\begin{bmatrix}
			76 & 2F & 61 & 2E\\
			66 & 28 & 60 & 29 \\
			50 & 18 & 53 & 61 \\
			D1 & 90 & D8 & EC
		\end{bmatrix}
		=
		\begin{bmatrix}
			C9 & E7 & AE & 49\\
			75 & 8B & D8 & D2 \\
			5D & 9F & CA & D9 \\
			51 & 14 & FC & F6
		\end{bmatrix}
		$
	\end{center}
	\item Thực hiện quá trình mã hóa với các bước tương tự như vậy cho đến vòng lặp cuối cùng. Ở vòng lặp cuối cùng, ta không làm bước \textbf{MixColumns}.
	
	
	\begin{table}[h]
		\centering
		\renewcommand{\arraystretch}{1.5}
		
		\setlength{\arraycolsep}{4pt} 
		
		\resizebox{\textwidth}{!}{
			\normalsize 
			\begin{tabular}{|c|c|c|c|c|c|c|}
				\hline
				\textbf{Round} & \textbf{Bắt đầu} & \textbf{Sau SubBytes} & \textbf{Sau ShiftRows} & \textbf{Sau MixColumns} & & \textbf{RoundKey} \\ \hline
				
				% Round 2
				Round 2 & 
				$\begin{array}{cccc} C9 & E7 & AE & 49 \\ 75 & 8B & D8 & D2 \\ 5D & 9F & CA & D9 \\ 51 & 14 & FC & F6 \end{array}$ & 
				$\begin{array}{cccc} DD & 94 & E4 & 3B \\ 9D & 3D & 61 & B5 \\ 4C & DB & 74 & 35 \\ D1 & FA & B0 & 42 \end{array}$ & 
				$\begin{array}{cccc} DD & 94 & E4 & 3B \\ 3D & 61 & B5 & 9D \\ 74 & 35 & 4C & DB \\ 42 & D1 & FA & B0 \end{array}$ & 
				$\begin{array}{cccc} D0 & 74 & A1 & A1 \\ 79 & D8 & BB & DC \\ CE & F7 & DC & C0 \\ B1 & 4A & 21 & 70 \end{array}$ &
				$\oplus$ &
				$\begin{array}{cccc} D1 & FE & 9F & B1 \\ 89 & A1 & C1 & E8 \\ 9E & 86 & D5 & B4 \\ E0 & 70 & A8 & 44 \end{array}$ \\ \hline
				
				% Hàng trống gộp cột
				\multicolumn{7}{|c|}{\dots} \\ \hline
				
				% Round 10
				Round 10 & 
				$\begin{array}{cccc} 8F & DF & A1 & 22 \\ 01 & 00 & 78 & 14 \\ B2 & 30 & E8 & 96 \\ AA & 1F & 2B & 16 \end{array}$ & 
				$\begin{array}{cccc} 73 & 9E & 32 & 93 \\ 7C & 63 & BC & FA \\ 37 & 04 & 9B & 90 \\ AC & C0 & F1 & 47 \end{array}$ & 
				$\begin{array}{cccc} 73 & 9E & 32 & 93 \\ 63 & BC & FA & 7C \\ 9B & 90 & 37 & 04 \\ 47 & AC & C0 & F1 \end{array}$ & 
				--- & 
				$\oplus$ &
				$\begin{array}{cccc} C8 & F2 & 79 & 54 \\ 76 & 97 & CE & 19 \\ 13 & 3E & 76 & E8 \\ 4C & 46 & 67 & 89 \end{array}$ \\ \hline
			\end{tabular}
		}
	\end{table}
	
\end{enumerate}

Output của vòng lặp cuối cùng, cũng chính là đầu ra của thuật toán mã hóa đối xứng AES:

\begin{center}
	$
	\text{AESEncryption(Plaintext, CypherKey)}
	=
	\begin{bmatrix}
		BB & 6C & 4B & C7\\
		15 & 2B & 34 & 65\\
		88 & AE & 41 & EC\\
		0B & EA & A7 & 78
	\end{bmatrix}
	$
	
\end{center}

Vậy \textbf{Ciphertext} tương ứng là: $BB15880B6C2BAEEA4B3441A7C765EC78$

